% Generated draft: requires mathematical and source review.
\documentclass[11pt]{article}
\usepackage[margin=25mm]{geometry}
\usepackage{fontspec}
\setmainfont{Latin Modern Roman}
\usepackage{amsmath,amssymb,mathtools,unicode-math}
\setmathfont{Latin Modern Math}
\usepackage{enumitem,needspace,xurl,hyperref}
\setlength{\parindent}{0pt}
\setlength{\parskip}{5pt}
\setlength{\emergencystretch}{3em}
\newcommand{\N}{\mathbb{N}}
\newcommand{\Z}{\mathbb{Z}}
\newcommand{\Q}{\mathbb{Q}}
\newcommand{\R}{\mathbb{R}}
\newcommand{\C}{\mathbb{C}}
\newcommand{\F}{\mathbb{F}}
\newcommand{\A}{\mathbb{A}}
\newcommand{\PP}{\mathbb{P}}
\newcommand{\E}{\mathbb{E}}
\newcommand{\eps}{\varepsilon}
\newcommand{\dd}{\,\mathrm{d}}
\begin{document}
\section*{UnsolvedMath ID: 701664}
\textbf{Source problem code:} catalogue entry
\subsection{Definitions and mathematical statement}
For a group \(X\), write \(\operatorname{ord}(x)\) for the order of an element \(x\in X\), and define
\[
\pi(X)=\{p\in\N:p\text{ is prime and }\operatorname{ord}(x)=p\text{ for some }x\in X\}.
\]
The prime graph, also called the Gruenberg--Kegel graph, \(\Gamma(X)\), is the simple undirected graph with vertex set \(\pi(X)\), in which distinct vertices \(p,q\) are adjacent precisely when \(X\) contains an element of order \(pq\). This definition also applies when \(X\) is infinite.

Let \(G\) be a finite group. Its integral group ring is
\[
\Z G=\left\{\sum_{g\in G}a_g g:a_g\in\Z\right\},
\]
with multiplication induced by the group multiplication and distributivity. The augmentation homomorphism is
\[
\varepsilon:\Z G\longrightarrow\Z,\qquad
\varepsilon\!\left(\sum_{g\in G}a_g g\right)=\sum_{g\in G}a_g.
\]
The group of normalized units is
\[
\mathrm{V}(\Z G)=\{u\in(\Z G)^\times:\varepsilon(u)=1\}.
\]

The question is whether, for every finite group \(G\),
\[
\Gamma(G)=\Gamma\bigl(\mathrm{V}(\Z G)\bigr)
\]
as graphs with vertices labeled by primes. Explicitly, the requested equality requires both
\[
\pi(G)=\pi\bigl(\mathrm{V}(\Z G)\bigr)
\]
and, for every pair of distinct primes \(p,q\),
\[
\bigl(\exists g\in G:\operatorname{ord}(g)=pq\bigr)
\quad\Longleftrightarrow\quad
\bigl(\exists u\in\mathrm{V}(\Z G):\operatorname{ord}(u)=pq\bigr).
\]

The supplied statement attributes the question to W. Kimmerle and reports that it has been reduced to almost simple groups. Here an almost simple group means a finite group \(H\) satisfying \(S\leq H\leq\operatorname{Aut}(S)\) for some nonabelian finite simple group \(S\), identified with its inner automorphism group. The reduction is supplied background attributed to W. Kimmerle and A. Konovalov; its precise theorem and hypotheses are not reproduced in the record and have not been independently verified here. The separate proposer metadata is malformed, so its intended attribution remains unresolved. The registry labels the problem open; that label is supplied source metadata.
\subsection{Short English statement}
Does every finite group have exactly the same prime graph as the group of augmentation-one units in its integral group ring? Equality includes the prime labels of the vertices and the existence of elements of order \(pq\) for every pair of distinct primes. The supplied source attributes the question to W. Kimmerle and reports a reduction to almost simple groups.
\subsection{Sources}
\begin{itemize}
\item Supplied local UnsolvedMath registry record, permanent ID \(701664\), exact source problem code \texttt{catalogue entry}, title \emph{Kourovka Notebook Problem 20.6}. The supplied registry status is \texttt{open}; this is not an independently verified literature assessment.
\item \emph{The Kourovka Notebook (Unsolved Problems in Group Theory)}, original problem number 20.6. The background identifies the 20th issue (2022) and the 21st issue edition (2026), with arXiv identifier \texttt{1401.0300v48}. The collection description names E. I. Khukhro and V. D. Mazurov as editors. These bibliographic details are supplied by the record; no source was visited.
\item The statement begins with the attribution W. Kimmerle. The separate proposer field is truncated after an incomplete accent encoding following A. B; the background repeats this corruption and places the fragment ``chle, L. Margolis'' with a closing brace beneath its discussion-and-literature heading. These fragments do not provide a reliable complete proposer attribution.
\item W. Kimmerle and A. Konovalov, \emph{Internat. J. Algebra Comput.}, volume 27 (2017), pages 619--631. The supplied statement cites this reference for the reduction to almost simple groups. No article title, theorem number, or URL is supplied, and the reduction has not been independently verified here.
\end{itemize}
\end{document}
